\documentclass[12pt]{article}
\usepackage[utf8]{inputenc}
\usepackage[T1]{fontenc}
\usepackage{hyperref}
\title{Kurumba, Jennu Language Resources for K-12 Curriculum Development}
\author{Compiled by Claude Sonnet 5.5 (claude-sonnet-5-5) for Andr\'as Kornai}
\date{}
\begin{document}
\maketitle

\section{Introduction}
Jennu (Jenu) Kurumba (ISO 639-3 xuj; Glottocode jenn1240) is a Southern Dravidian language. It is often treated as a Kannada dialect, but Ethnologue lists it separately. Speakers live around the Nilgiri Hills border area of Tamil Nadu, Karnataka and Kerala (Mysore and Kodagu districts, and Wayanad). Estimates differ: Wikipedia gives about 100,000 native speakers (2011 census), while UNESCO's atlas lists 36,076. The community is the Jenu Kuruba, traditionally honey gatherers. I found no school materials in the language.

\section{Primary Educational Resources}
I found no primers, textbooks or school programmes for Jennu Kurumba in the search results.

\section{Linguistic and Reference Materials}
\begin{itemize}
\item Jenu Kurumba language (Wikipedia). \url{https://en.wikipedia.org/wiki/Jenu_Kurumba_language}
\item Jenu Kuruba, an article on the people. \url{https://en.wikipedia.org/wiki/Jenu_Kuruba}
\item UNESCO Atlas of the World's Languages in Danger entry. It lists only a few descriptions, mostly from the 1970s and 1980s. \url{https://ar.wal.unesco.org/index.php/countries/india/languages/jennu-kurumba}
\item Grambank entry (jenn1240). \url{https://grambank.clld.org/languages/jenn1240}
\item CIIL Sanchika repository item. I did not open it, so its contents are unverified. \url{https://sanchika.ciil.org/handle/1/11916}
\end{itemize}

\section{Religious and Community Resources}
\begin{itemize}
\item Oxford University Press chapter on identity politics and ethical worlds in South India. It appeared in the search results, and I did not verify that it treats Jenu Kuruba specifically. \url{https://academic.oup.com/book/27193/chapter/196635937}
\end{itemize}

\section{Digital and Audio Resources}
\begin{itemize}
\item DoBeS (MPI Nijmegen) documentation project page. It reports a documentation project on the language. \url{https://dobes.mpi.nl/?p=1168}
\item DoBeS poster on Kurumba. \url{https://dobes.mpi.nl/posters/print/eurasia-kurumba-print.pdf}
\end{itemize}

\section{Curriculum Development Context}
No school is known to teach Jennu Kurumba, and children are generally schooled in Kannada, Tamil or Malayalam. The DoBeS documentation may offer recordings and texts to build on. Community participation is needed, and local educational bodies and tribal welfare departments are possible partners, but I found none named in the search results.

\end{document}